\scanpage{89}
% archive foliation: 88
De (vi) et (viii) résulte que pour tt
motif $M$ et entier $n$, l'isomorphisme
\[
\Phi^s(M)\simeq\Phi^s(M(n))
\]
est compatible avec les cônes amples.

De (iv) et \unsure{(vi)|(vii)} résulte que si $\rho=0$, une
forme ample sur $M$ induit une
forme \uline{définie positive}
\[
\Gamma_k M\times\Gamma_k M\longrightarrow\Gamma_k\mathbf Q(0)=\mathbf Q
\]

Si on décompose $M$ (de poids $\rho$) en
composantes isotypiques $M_\alpha$, il vient
\[
\Phi^s(M)\simeq\unsure{\prod|\bigoplus}_\alpha\Phi^s(M_\alpha)
\]
donc pour étudier les formes positives, il
suffit de les étudier en principe sur
des morceaux isotypiques, qui sont
d'ailleurs \unsure{sous|de} la forme
\[
I\otimes_A\unsure{V|W}
\]
où $I$ est un motif simple, $A$ son
anneau des endomorphismes, qui est
un corps, et $\unsure{V|W}$ un \uline{vectoriel} sur $A$
